\section{Additional Security and Accountability Measurements}
\label{app:e-security-accountability}

\noindent\textbf{Current valid-signature adaptive suite.}
Table~\ref{tab:e-attack-families} expands the current protocol attack results
by attack family. Each scheme has 100 episodes in total, covering six families;
there are not 100 episodes per family. Each episode permits at most 16 queries
and stops at the first accepted violation. The next attack family can depend on
the first-failed-relation feedback exposed by the auditor. Every retained attack
query passed both a signature-validity check and an independent protocol-violation
oracle. Family counts therefore differ across schemes. SignedLog+ and
\textsc{TokenProv} both reject all tested queries; these observations do not
establish a detection-rate advantage of \textsc{TokenProv} over SignedLog+.
These are CPU protocol microbenchmarks, rather than attacks executed against
live model inference.

\begin{table*}[!t]
\centering
\caption{Current valid-signature adaptive attacks by family (accepted attack queries / actual queries). Family query counts differ because search adapts to feedback and stops at the first success.}
\label{tab:e-attack-families}
\footnotesize
\renewcommand{\arraystretch}{1.10}
\begin{tabular*}{\textwidth}{@{\hspace{4pt}}l@{\extracolsep{\fill}}cccc@{\hspace{4pt}}}
\toprule
\textbf{Attack family} & \textbf{SignedLog} & \textbf{SignedLog+} & \textbf{TokenProv} & \shortstack{\textbf{TokenProv first}\\\textbf{failed relation}} \\
\midrule
Aggregate bill manipulation & 18/18 & 0/264 & 0/232 & $R_{\mathrm{set}}$ \\
Receipt duplication & 13/13 & 0/267 & 0/223 & $R_{\mathrm{prov}}$ \\
Receipt omission & 22/22 & 0/278 & 0/256 & $R_{\mathrm{set}}$ \\
Observable accepted-call omission & 12/12 & 0/265 & 0/283 & $R_{\mathrm{hand}}$ \\
Cross-call evidence splicing & 14/14 & 0/251 & 0/309 & $R_{\mathrm{prov}}$ \\
Unauthorized deployment substitution & 21/21 & 0/275 & 0/297 & $R_{\mathrm{prov}}$ \\
\bottomrule
\end{tabular*}
\end{table*}


\noindent\textbf{Historical protocol fixtures.}
Tables~\ref{tab:e-historical-attack-coverage}--\ref{tab:e-historical-collusion}
report retained historical microbenchmarks separately from the current suite.
Each attack, scenario--variant pair, or injected stage has 500 trials across
five fixture seeds, with 100 repetitions per seed. These repetitions do not
represent 500 distinct attacks or independent serving deployments.
Table~\ref{tab:e-historical-attack-coverage} includes attacks that alter signed
fields without repairing their original signatures; its detection results
must not be pooled with the current valid-signature adaptive suite.

\begin{table*}[!t]
\centering
\caption{Historical protocol attack microbenchmark (detected / 500 trials for each attack and scheme). These fixtures are separate from the current valid-signature adaptive suite.}
\label{tab:e-historical-attack-coverage}
\footnotesize
\renewcommand{\arraystretch}{1.10}
\begin{tabular*}{\textwidth}{@{\hspace{4pt}}l@{\extracolsep{\fill}}lccc@{\hspace{4pt}}}
\toprule
\textbf{Attack} & \textbf{Expected relation} & \textbf{Vanilla} & \textbf{SignedLog} & \textbf{TokenProv} \\
\midrule
A1: Unauthorized call & $R_{\mathrm{auth}}$ & 0/500 & 500/500 & 500/500 \\
A2: Call insertion & $R_{\mathrm{life}}$ & 0/500 & 0/500 & 500/500 \\
A3: Observable-call omission & $R_{\mathrm{hand}}$ & 0/500 & 0/500 & 500/500 \\
A4: Call replay & $R_{\mathrm{prov}}$ & 0/500 & 0/500 & 500/500 \\
A5: Call relabeling & $R_{\mathrm{prov}}$ & 0/500 & 500/500 & 500/500 \\
A6: $H\!\rightarrow\!B$ boundary tampering & $R_{\mathrm{hand}}$ & 0/500 & 500/500 & 500/500 \\
A7: $B\!\rightarrow\!P$ boundary tampering & $R_{\mathrm{hand}}$ & 0/500 & 500/500 & 500/500 \\
A8: Model/deployment substitution & $R_{\mathrm{prov}}$ & 0/500 & 500/500 & 500/500 \\
A9: Tokenizer/input-semantics substitution & $R_{\mathrm{prov}}$ & 0/500 & 500/500 & 500/500 \\
A10: Cache or usage inflation & $R_{\mathrm{prov}}$ & 0/500 & 500/500 & 500/500 \\
A11: Wrong price version & $R_{\mathrm{prov}}$ & 0/500 & 500/500 & 500/500 \\
A12: Cross-call evidence splicing & $R_{\mathrm{prov}}$ & 0/500 & 0/500 & 500/500 \\
A13: Receipt omission or duplication & $R_{\mathrm{set}}$ & 0/500 & 0/500 & 500/500 \\
A14: Aggregate bill manipulation & $R_{\mathrm{set}}$ & 0/500 & 500/500 & 500/500 \\
\bottomrule
\end{tabular*}
\end{table*}


\noindent\textbf{Cross-call evidence binding.}
Table~\ref{tab:e-historical-splicing} covers same-deployment, cross-deployment,
and cross-model source--target pairs. The source and target requests differ,
while the fixture call label can remain \texttt{c-0}.
IndependentLocalChecks verifies the original receipt under the source
provider key and registration context; the full auditor checks the transplanted
provider evidence against the target request transcript.
The archived \texttt{TokenProv-w/o-R\_link} variant executes the same receipt-only
local-validity branch as IndependentLocalChecks. We therefore report it as a
local-check surrogate, rather than as a strict single-check ablation of the
complete auditor. Each of the three variants has 1,500 trials in total;
the archived file contains 4,500 rows across all variants.

\begin{table*}[!t]
\centering
\caption{Historical cross-call splicing microbenchmark (accepted / 500 trials). The local-check surrogate executes the same receipt-only check as IndependentLocalChecks; it is not a strict one-check removal from the complete auditor.}
\label{tab:e-historical-splicing}
\footnotesize
\renewcommand{\arraystretch}{1.10}
\begin{tabular*}{\textwidth}{@{\hspace{4pt}}l@{\extracolsep{\fill}}ccc@{\hspace{4pt}}}
\toprule
\textbf{Scenario} & \shortstack{\textbf{Independent}\\\textbf{local checks}} & \shortstack{\textbf{Local-check}\\\textbf{surrogate}} & \textbf{TokenProv} \\
\midrule
S1: Same model and deployment; different request & 500/500 & 500/500 & 0/500 \\
S2: Same model; different deployment and request & 500/500 & 500/500 & 0/500 \\
S3: Different model, deployment, and request & 500/500 & 500/500 & 0/500 \\
\bottomrule
\end{tabular*}
\end{table*}


\noindent\textbf{First-failure localization.}
Table~\ref{tab:e-historical-localization} reports agreement between the first
failed relation and the predefined injected stage, together with the returned
responsibility intervals and their widths. The source field named
\texttt{ground\_truth\_contained} duplicates the first-relation agreement test;
it is not an independent interval-membership check and is not reported as a
separate containment result.
Interval width counts hops along the $U\!\rightarrow\!H\!\rightarrow\!B\!\rightarrow\!P$
role path; zero width denotes a single role. The five stage-specific mutations
are predefined protocol fixtures. Agreement in these fixtures does not imply
unique malicious-party identification for arbitrary or simultaneous faults.

\begin{table*}[!t]
\centering
\caption{Historical five-stage fault-localization microbenchmark. Matching means the returned first failed relation equals the predefined injected-stage relation; the returned responsibility interval is fixed by that relation, and its width counts role-boundary hops.}
\label{tab:e-historical-localization}
\footnotesize
\renewcommand{\arraystretch}{1.10}
\begin{tabular*}{\textwidth}{@{\hspace{4pt}}l@{\extracolsep{\fill}}lcccc@{\hspace{4pt}}}
\toprule
\textbf{Injected stage} & \shortstack{\textbf{First failed}\\\textbf{relation}} & \shortstack{\textbf{Relation}\\\textbf{matches}} & \shortstack{\textbf{Returned}\\\textbf{interval}} & \shortstack{\textbf{Mean interval}\\\textbf{width (hops)}} & \shortstack{\textbf{Mean audit}\\\textbf{time (ms)}} \\
\midrule
Authorization & $R_{\mathrm{auth}}$ & 500/500 & $U$--$H$ & 1 & 0.168 \\
Lifecycle & $R_{\mathrm{life}}$ & 500/500 & $H$--$H$ & 0 & 3.476 \\
Handoff & $R_{\mathrm{hand}}$ & 500/500 & $H$--$P$ & 2 & 3.430 \\
Provider & $R_{\mathrm{prov}}$ & 500/500 & $B$--$P$ & 1 & 9.165 \\
Settlement & $R_{\mathrm{set}}$ & 500/500 & $H$--$B$ & 1 & 14.350 \\
\bottomrule
\end{tabular*}
\end{table*}


\noindent\textbf{Collusion and evidence observability.}
Table~\ref{tab:e-historical-collusion} distinguishes detectable contradictions
from explicit evidence-boundary controls. A retained signed boundary exposes
the $B+P$ splicing case. In the $H+B$ hidden-execution and internally consistent
$H+B+P$ controls, no conflicting independent evidence is supplied, and the
auditor accepts the available transcript. These accepted controls are retained
and are not counted as blocked attacks. The measurements do not test
evidence-expiry or non-submission enforcement.

\begin{table*}[!t]
\centering
\caption{Historical collusion and observable-evidence controls. Acceptance in the last two rows is an evidence-boundary control, not a failed rejection of an observably inconsistent transcript.}
\label{tab:e-historical-collusion}
\footnotesize
\renewcommand{\arraystretch}{1.10}
\begin{tabular*}{\textwidth}{@{\hspace{4pt}}l@{\extracolsep{\fill}}p{.20\textwidth}p{.25\textwidth}p{.14\textwidth}p{.12\textwidth}@{\hspace{4pt}}}
\toprule
\textbf{Corrupt roles} & \textbf{Scenario} & \textbf{Independent evidence} & \textbf{Outcome} & \shortstack{\textbf{Responsibility}\\\textbf{interval}} \\
\midrule
$B+P$ & Cross-call provider-evidence splicing & Retained signed $H\!\rightarrow\!B$ boundary conflicts with provider-linked evidence & 500/500 rejected & $B$--$P$ \\
$H+B$ & One hidden execution omitted from an otherwise consistent transcript & No independent receipt for the hidden execution & 500/500 accepted & Unavailable \\
$H+B+P$ & Internally consistent service-side transcript & No conflicting external evidence anchor & 500/500 accepted & Unavailable \\
\bottomrule
\end{tabular*}
\end{table*}

